\chapter*{Start here: from zero}
\addcontentsline{toc}{chapter}{Start here: from zero}
\markboth{Start here}{Start here}

Physics is the study of how things move and interact, written in mathematics.
This book covers the physics an engineer uses: motion, forces, energy,
rotation, electricity, magnetism, waves and heat. You need algebra, a little
trigonometry (sine and cosine of an angle), and three tools that this chapter
teaches from scratch: units, a first look at calculus, and vectors.

\section*{Units: the first check on every answer}

Every physical quantity has a unit, and a formula must balance in units. A
distance is in metres (m), a time in seconds (s), a speed in m/s, an acceleration
in m/s$^2$, a force in newtons (N, where $1\,\mathrm{N}=1\,\mathrm{kg\,m/s^2}$).
Check a formula by its units: the distance travelled from rest with acceleration
$a$ is $x=\tfrac12at^2$. In units, $(\mathrm{m/s^2})\cdot\mathrm{s^2}=\mathrm{m}$, which is a distance, so the formula is
at least plausible; a claimed $x=at$ would give m/s, a speed, so it cannot be right.
This habit catches more errors than any other, and every chapter uses it.

\section*{Calculus in two ideas}

\textbf{The derivative is a slope.} If a car's position is $x(t)=5t^{2}$ metres at
time $t$ seconds, its position at $t=2$ is $20$ m and at $t=2.01$ is $20.2005$ m.
Over that $0.01$ s it moved $0.2005$ m, a speed of about $20.05$ m/s. Shrink the
interval to nothing and you get exactly $20$ m/s. The rule that produces it is the
\emph{derivative}: for $x(t)=5t^{2}$, the velocity is $v(t)=\mathrm{d}x/\mathrm{d}t=10t$, and $v(2)=20$ m/s.
The derivative of $t^{n}$ is $nt^{n-1}$; that is the only rule Chapters~1--4 need.

\textbf{The integral is an accumulated total.} If the acceleration is a constant
$3\,\mathrm{m/s^2}$ for $4$ s, the velocity gained is $3\times4=12$ m/s: the area of
a rectangle of height $3$ and width $4$ on an acceleration--time graph. If the acceleration
varies, you add up many thin rectangles; that sum is the \emph{integral},
$\Delta v=\int a\,\mathrm{d}t$. The integral is the derivative run backwards:
since $v=10t$ is the derivative of $5t^{2}$, the integral of $10t$ from $0$ to $2$
is $5\cdot2^{2}-0=20$.

Chapter~1 uses both ideas to derive the constant-acceleration formulas from
definitions, so you do not memorise them.

\section*{Vectors: quantities with a direction}

A \emph{scalar} has size only (a mass, a time). A \emph{vector} has size and
direction (a velocity, a force). In a plane we write a vector with
\emph{components}: $\mathbf v=(3,4)$ means $3$ along $x$ and $4$ along $y$. Its
length is $|\mathbf v|=\sqrt{3^{2}+4^{2}}=5$. Adding vectors adds components:
$(3,4)+(1,-2)=(4,2)$.

Two ways of multiplying vectors appear later in the book:
\begin{description}
\item[Dot product] $\mathbf a\cdot\mathbf b=a_xb_x+a_yb_y+a_zb_z$, a number. It
      is $|\mathbf a||\mathbf b|\cos\theta$, so it is zero when the vectors are
      perpendicular. Work is a dot product (Chapter~4): $(1,2)\cdot(3,4)=3+8=11$.
\item[Cross product] $\mathbf a\times\mathbf b$, a vector perpendicular to both,
      with components $(a_yb_z-a_zb_y,\ a_zb_x-a_xb_z,\ a_xb_y-a_yb_x)$. Torque and
      the magnetic force are cross products (Chapters~5, 6 and~13). For
      $(1,2,3)\times(4,5,6)=(-3,\,6,\,-3)$, and $\hat{\mathbf x}\times\hat{\mathbf y}=\hat{\mathbf z}$.
\end{description}
You can check the cross-product example: $2\cdot6-3\cdot5=-3$, $3\cdot4-1\cdot6=6$, $1\cdot5-2\cdot4=-3$.

\section*{How the books of this series fit together}

\begin{center}
\begin{tabular}{P{2.4cm}P{9.0cm}}
\toprule
\thead{Book} & \thead{Connection}\\
\midrule
5 Electronics & Applies the circuit and field ideas of Chapters~9--14 here in its capacitor, inductor and power chapters. Each chapter here ends with an ``In engineering practice'' section that points to Book~5.\\
3 Linear algebra & Gives the vector and matrix machinery in greater depth; Chapters~1--5 there are a natural companion to the vector work here.\\
\bottomrule
\end{tabular}
\end{center}

\section*{How to use the boxes in each chapter}

Each chapter opens with \emph{After this chapter you can}. From Chapter~2, a
\emph{Before you start} box lists what the chapter relies on, with a recap.
Always sanity-check your answer with units and with an extreme case (what if the
mass is zero, or the angle is $90^\circ$?) before you trust it.
